% !TEX root = ./main.tex
\chapter*{符号与记号}
\addcontentsline{toc}{chapter}{符号与记号}

本书采用时间因子 $\ee^{-\ii\omega t}$ 约定。除非特别说明，$\omega$ 为角频率，$f=\omega/(2\pi)$，波长为 $\lambda$。

\begin{longtable}{p{0.20\textwidth}p{0.72\textwidth}}
\toprule
符号 & 含义 \\
\midrule
\endhead
$\rvec=(x,y,z)$ & 直角坐标；$z$ 轴取外延层法向。 \\
$\kvec$ & 波矢；在周期系统中写为 $\kvec+\gvec$。 \\
$\gvec$ & 倒格矢（reciprocal lattice vector）。 \\
$\Gamma$ 点 & 布里渊区中心，$\kvec=\bm{0}$。 \\
$\E,\Hf$ & 频域电场（electric field）与磁场（magnetic field）复振幅。 \\
$\varepsilon(\rvec,\omega)$ & 介电函数；可含色散与损耗。 \\
$Q$ & 品质因子，$Q=\omega_0/(2\gamma)$（被动定义）。 \\
$\gamma_\mathrm{rad},\gamma_\mathrm{int}$ & 辐射损耗与内部损耗对应的衰减率。 \\
$\kappa_1,\kappa_2$ & CWT 中由主导几何傅里叶分量和纵向模重叠给出的耦合系数。 \\
$\betaslab$ & slab/导波模型中的面内传播常数（单位：\si{\per\meter}）。 \\
$\betasp$ & 自发辐射耦合因子（无量纲），定义见动态与噪声章节。 \\
$\betabloch$ & 仅在少数上下文中用于记 Bloch 波矢分量；默认更推荐直接写 $\kvec$ 分量以避免歧义。 \\
$g_\mathrm{mat}$ & 材料增益（material gain）。 \\
$g_N$ & 微分增益（differential gain），即材料增益对载流子密度的局域导数。 \\
$N_\mathrm{tr}$ & 透明载流子密度（transparency carrier density）。 \\
$\Gamma_\mathrm{opt}$ & 光学约束因子（optical confinement factor）。 \\
$\Delta G_{\star j}$ & 模鉴别裕量；表示目标模与第 $j$ 个竞争模之间的净优势。 \\
$g_\mathrm{mod}$ & 模增益（modal gain），常写为 $\Gamma_\mathrm{opt}g_\mathrm{mat}$（近似）。 \\
$g_\mathrm{th}$ & 阈值增益（threshold gain），属于光学阈值条件。 \\
$I_\mathrm{th}$ & 阈值电流（threshold current），属于器件级电-热-光结果。 \\
$N$ & 载流子浓度（carrier density）。 \\
$\phi$ & 静电势（electrostatic potential）。 \\
$T$ & 温度场。 \\
$R_\square$ & 电流扩展层的片电阻（sheet resistance）。 \\
$L_s$ & 电流扩展长度（spreading length）。 \\
\bottomrule
\end{longtable}

\begin{PitfallBox}
符号 $Q$、$g_\mathrm{th}$、$I_\mathrm{th}$ 分别对应被动腔、光学阈值、器件阈值三个层级，不能互相替代。$\beta$ 统一优先写成 $\betaslab$ 或 $\betasp$，避免跨章节误读。
\end{PitfallBox}

\chapter*{缩略语表}
\addcontentsline{toc}{chapter}{缩略语表}

\begin{longtable}{p{0.20\textwidth}p{0.72\textwidth}}
\toprule
缩写 & 含义 \\
\midrule
\endhead
PCSEL & photonic-crystal surface-emitting laser，光子晶体面发射激光器。 \\
PWEM & plane-wave expansion method，平面波展开法。 \\
RCWA & rigorous coupled-wave analysis，严格耦合波分析。 \\
FDTD & finite-difference time-domain，时域有限差分法。 \\
FEM & finite-element method，有限元法。 \\
CWT & coupled-wave theory，耦合波理论。 \\
TCMT & temporal coupled-mode theory，时间耦合模理论。 \\
BIC & bound state in the continuum，连续谱中的束缚态。 \\
QNM & quasinormal mode，准正规模。 \\
SCH & separate confinement heterostructure，分离限制异质结构。 \\
SQW / MQW & single quantum well / multiple quantum wells，单量子阱 / 多量子阱。 \\
PML & perfectly matched layer，完全匹配层。 \\
EBL / HBL & electron blocking layer / hole blocking layer，电子阻挡层 / 空穴阻挡层。 \\
RIN & relative intensity noise，相对强度噪声。 \\
SALT & steady-state \emph{ab initio} laser theory，稳态从头算激光理论。 \\
\bottomrule
\end{longtable}
